\scanpage{36}
\mspage{17}
and this being posited, relation (56) is equivalent to
\[
(60)\qquad\left\{\begin{aligned}
\alpha'&=\zeta'\rho^{\deleted{\unsure{-1}}}(\zeta')^{-1}\\
\text{i.e.}\quad\eta\sigma_\infty(\eta)^{-1}&=(\zeta'\rho^{\deleted{\unsure{-1}}}(\zeta')^{-1})\sigma_\infty(\lambda)
\end{aligned}\right.
\]
\deleted{i.e. to}\quad\deleted{(61)}
Indeed, formula (75bis) expressing $\alpha$ in terms of $\widetilde\alpha$
\[
(x)\qquad\alpha=\rho(\eta\widetilde\alpha\rho^{-1}(\eta)^{-1})
\]
is rewritten, on taking inverses and setting
\[
(xx)\qquad\widetilde\alpha'=\widetilde\alpha^{-1}
\]
in the form
\[
(61)\qquad\uline{\alpha'}=\rho(\rho^{-1}(\eta)\widetilde\alpha'\eta^{-1})
=\uline{\eta\rho(\widetilde\alpha')\rho(\eta)^{-1}}
\]
\deleted{and the formula}
(Here we consider $\alpha,\alpha',\widetilde\alpha,\widetilde\alpha'$ as deduced from $\lambda,\eta\in\widehat\Pi_{0,3}$ by the first relation (55) as regards $\widetilde\alpha$, from which the others are deduced by $(x),(xx)$. Moreover, the \inserted{2\textsuperscript{nd}} formula (55) is\deleted{\unsure{t}} then rewritten, on taking inverses
\[
\widetilde\alpha'=\rho^{-1}(\zeta)\zeta^{-1}
\]
i.e.
\[
\rho(\widetilde\alpha')=\zeta\rho(\zeta)^{-1}
\]
\deleted{so that (61), \unsure{(62)} can also be written}
\inserted{which is equivalent} \inserted{\unsure{(61)}}
\[
\alpha'=\eta\zeta\rho(\zeta^{-1}\eta^{-1})
\]
\deleted{i.e. $\alpha'$ to (60)},
i.e. to (60), with
\[
\zeta'=\eta\zeta.
\]

\uline{In summary} The automorphisms of the groupoid $\overline{\underline{\mathrm{VI}}}_J$, inducing the identity on the vertices and \uline{commuting with the action of $G_J$}, correspond to the
\[
(62)\qquad\boxed{(\lambda,\eta,\zeta')\in\widehat\Pi_{0,3}\times\widehat\Pi_{0,3}\times\widehat\Pi_{0,3}}
\]
satisfying the loop relation ?
\margin{35}
